\section{Design}
\label{sec:design}

% Target: ~2 two-column pages (~4 single-column equivalent)
% This is the core contribution section. Condense report sections/analysis.tex.
% Present the vulnerability analysis, solution-space exploration, and the
% chosen domain-scheduler design.

\subsection{FreeRTOS Scheduler Vulnerability}
\label{sub:vulnerability}

% [~0.5 column -- motivate the need for redesign]
%   - FreeRTOS uses fixed-priority preemptive scheduling with optional
%     round-robin among equal-priority tasks.
%   - The scheduler is inherently unfair by design: a higher-priority task
%     that never yields starves all lower-priority tasks.
%   - Timing-channel exploit: a high-priority Trojan encodes a secret bit
%     by choosing to sleep (allowing lower-priority spy to run immediately)
%     or spin (delaying spy). The spy measures its own wake-up latency via
%     the system tick counter to infer the bit.
%   - This is not a bug but a consequence of priority semantics. The channel
%     exists across all priority levels unless time is decoupled from
%     scheduling priority.

\subsection{Approaches to Closing Scheduling-Based Channels}
\label{sub:approaches}

% [~1 column -- evaluate alternatives, justify choice]
%   - Approach 1: Remove time access. If tasks cannot read the system clock
%     or tick counter, timing channels through the scheduler vanish. Drawback:
%     real-time systems \emph{require} time for correct operation (deadlines,
%     periodic tasks). Infeasible for FreeRTOS.
%   - Approach 2: Protection within priority. Partition a single priority
%     level into fixed execution quanta, ensuring that two tasks at the same
%     priority cannot observe each other's execution time. Drawback: does
%     not address cross-priority channels; higher-priority tasks still leak.
%   - Approach 3: Domain scheduling. Introduce an abstraction layer above
%     priorities: groups of tasks (domains) that are collectively allocated
%     fixed time slices in a round-robin fashion. Within each domain, the
%     existing priority scheduler runs unchanged. Only domain-boundary
%     switches are externally observable.
%   - Chosen approach: domain scheduling (Approach 3). It provides the
%     strongest isolation with minimal changes to the existing scheduler.

\subsection{Domain-Level Scheduler Design}
\label{sub:domain_design}

% [~1.5 columns -- detailed design description]
%   - Domain 0 is the initial domain owning all time slices. Each subsequent
%     domain creation carves a contiguous range of slices from Domain 0.
%   - Time slices are numbered 0 through configNUM\_TIME\_SLICES $-$ 1.
%     configNUM\_TICKS\_PER\_SLICE determines granularity.
%   - At boot, Domain 0 has all slices. When Domain 1 is created with
%     (start\_slice = 4, num\_slices = 3), slices 4--6 are transferred.
%     Domain 0 retains slices 0--3 and 7--(N$-$1).
%   - The domain scheduler iterates through time slices in order. On reaching
%     a slice owned by a different domain, it triggers a domain switch.
%   - Domain switch procedure:
%     \begin{enumerate}
%       \item Save current domain's task context (existing FreeRTOS
%         mechanism).
%       \item Execute fence.t to flush microarchitectural state.
%       \item Restore next domain's ready task.
%     \end{enumerate}
%   - Within a domain, the standard priority preemptive scheduler operates
%     unchanged. Tasks are only preempted by higher-priority tasks \emph{in
%     the same domain}.

\subsection{Constant-Time Domain Creation}
\label{sub:constant_time_creation}

% [~0.75 column]
%   - Naive approach: scan the time-slice array to find free slices.
%     Information leak: the scan time reveals how many domains already exist
%     (i.e., the amount of free space).
%   - Design: an array `xDomains[configNUM\_TIME\_SLICES]` indexed by slice
%     number, combined with a `uxNext` singly-linked list through Domain 0's
%     unowned slices.
%   - Domain 0 maintains a free-list head pointer. Creating a domain:
%     \begin{enumerate}
%       \item Walk `num\_slices` steps in the free list (constant work per
%         slice irrespective of total system load).
%       \item Assign the traversed slice entries to the new domain.
%       \item Splice the remainder of the free list past the removed region.
%     \end{enumerate}
%   - Split operations on the left/right fragments are O(1) using the list
%     structure. No dependence on the number of existing domains.

% [Optional figure: Domain timeline diagram showing slice ownership]
%   - Include a TikZ or imported figure from report/figures/ showing the
%     time-slice array with colour-coded domain ownership.

\subsection{Cross-Domain Communication}
\label{sub:cross_domain_comm}

% [~0.5 column]
%   - FreeRTOS queues move tasks between ready lists when data is sent/
%     received. This can move a task across domain boundaries, introducing
%     a timing channel constrained to tasks sharing a queue.
%   - Design decision: accept this residual channel. Each TCB tracks its
%     `uxDomainID`. When a task is unblocked by a queue operation, it
%     resumes in its home domain's next available time slice.
%   - Alternative considered: prohibit cross-domain queues entirely. Rejected
%     because inter-domain communication is essential for practical systems.
%   - A synchronisation primitive (e.g., lock) shared across domains remains
%     an acknowledged timing channel, discussed further in Section VII.
